\begin{frame}{Self-Attention: 문장이 자기 자신을 본다}
  Q, K, V가 모두 \textbf{같은 문장}에서 나오면 \kb{self-attention}
  — 각 단어가 같은 문장 안의 다른 모든 단어(\textbf{자기 자신
  포함})와의 관련도를 계산. \textbf{이게 이번 장 도입부의
  메커니즘}: ``그것은''의 Query가 문장 안 모든 단어의 Key와
  내적했을 때 ``그 동물''의 Key와 \textbf{가장 높은 점수}가
  나오도록 학습된다.
  \vskip0.4em
  \begin{block}{구조를 한 줄로}
    각 단어의 ``새로운 표현'' = 같은 문장 안의 \emph{모든}
    단어(자기 자신 포함)의 Value 벡터를, ``그 단어가 각각과
    얼마나 관련 있는가''를 softmax로 나눠 \textbf{가중평균한 값}.
    \vskip0.3em
    \textbf{각 단어의 출력이 서로 \emph{독립}이 \emph{아니다}} —
    앞 네 단어 예에서 ``동물''의 출력에 ``그것은''의 Value가
    41.2\% 섞여 있는 \emph{동시에}, ``그것은''의 출력에도
    ``동물''의 Value가 56.1\% 섞인다. 두 단어는 서로의 표현을
    \emph{서로에게} 공급하면서 \textbf{동시에 업데이트} —
    RNN의 \emph{순차적} 은닉 상태 업데이트와 \textbf{가장
    근본적으로 다른 점} (12.1절의 병렬성 문제 바로 그 차이).
  \end{block}
  \vskip0.4em
  또 하나: \textbf{어텐션 연산 자체에는 ``단어의 위치''라는
  정보가 전혀 없다.} $QK^T$ 는 \emph{내적}만 계산 — 단어를
  \emph{재배열}해도 각 쌍의 원점수는 그대로(\textbf{순서 불변}),
  ``강아지가 고양이를 쫓는다''와 ``고양이가 강아지를 쫓는다''가
  \emph{같은} 출력을 내는 것은 이 순서 불변성의 \emph{직접 결과}.
  해결책 \kb{positional encoding}(12.3절) — 여기서는
  ``\textbf{어텐션이 \emph{순서를 모른다}는 사실 자체}''를
  기억해 두자.
\end{frame}
